\documentclass[12pt,a4paper]{book}

\input{../1ere-preambule}

\newcommand{\va}[1]{\lvert \ #1 \ \rvert}
\newcommand{\vaf}[1]{\left\lvert \ #1\  \right\rvert}

\begin{document}
\graphicspath{{images/}}
\setcounter{chapter}{2}
\chapter[Correction des exercices]{Produit Scalaire \\ correction des exercices}
\thispagestyle{fancy}

\section*{Définition du produit scalaire}

\subsection*{Ex 23 p 243}

\begin{enumMet}
    \item $\mathrm{AB}=1, \mathrm{AC}=5$ et $(\vv{\mathrm{AB}}, \vv{\mathrm{AC}})=40^{\circ}$\\
    $\vv{AB}\cdot\vv{AC}=AB\times AC\times\cos(\vv{AB},\vv{AC})$\\
    $\vv{AB}\cdot\vv{AC}\approx1\times5\times0,77\approx3,83$
    \item $\mathrm{AB}=2,5, \mathrm{AC}=3$ et $(\vv{\mathrm{AB}}, \vv{\mathrm{AC}})=70^{\circ}$\\
    $\vv{AB}\cdot\vv{AC}=AB\times AC\times\cos(\vv{AB},\vv{AC})\approx2,5\times3\times0,34\approx2,56$
    \item $\mathrm{AB}=4, \mathrm{AC}=\dfrac{5}{3}$ et $(\vv{\mathrm{AB}}, \vv{\mathrm{AC}})=-110^{\circ}$\\
    $\vv{AB}\cdot\vv{AC}=AB\times AC\times\cos(\vv{AB},\vv{AC})\approx4\times\dfrac{5}{3}\times(-0,34)\approx-2,28$
\end{enumMet}

\section*{Propriétés du produit scalaire}

\subsection*{Ex 24 p 243}

\begin{enumMet}
    \item $\mathrm{AB}=2, \mathrm{AC}=6$ et $(\vv{\mathrm{AB}}, \vv{\mathrm{AC}})=\dfrac{\pi}{4}$\\
    $\vv{AB}\cdot\vv{AC}=AB\times AC\times\cos(\vv{AB},\vv{AC})=2\times6\times\dfrac{\sqrt{2}}{2}=6\sqrt{2}$
    \item $\mathrm{AB}=\sqrt{3}, \mathrm{AC}=1$ et $(\vv{\mathrm{AB}}, \vv{\mathrm{AC}})=-\dfrac{\pi}{6}$\\
    $\vv{AB}\cdot\vv{AC}=AB\times AC\times\cos(\vv{AB},\vv{AC})=\sqrt{3}\times1\times\dfrac{\sqrt{3}}{2}=\dfrac{3}{2}$
    \item $\mathrm{AB}=5, \mathrm{AC}=\dfrac{3}{4}$ et $(\vv{\mathrm{AB}}, \vv{\mathrm{AC}})=\dfrac{2 \pi}{3}$\\
    $\vv{AB}\cdot\vv{AC}=AB\times AC\times\cos(\vv{AB},\vv{AC})=5\times\dfrac{3}{4}\times\left(-\dfrac{1}{2}\right)=-\dfrac{15}{8}$
\end{enumMet}

\subsection*{Ex 51 p 246}

\begin{enumMet}
    \item
    \item
    \item $( \vv{u}- \vv{v}) \cdot( \vv{u}+ \vv{v}) = \vv{u} \cdot  \vv{u}- \vv{v} \cdot  \vv{v} =\| \vv{u}\|^{2}-\| \vv{v}\|^{2}=6^{2}-5^{2}=36-25=11$
    \item \ \\
    $\begin{aligned}( \vv{u}+2  \vv{v}) \cdot(3  \vv{u}+ \vv{v}) &= \vv{u} \cdot 3  \vv{u}+ \vv{u} \cdot  \vv{v}+2  \vv{v} \cdot 3  \vv{u}+2  \vv{v} \cdot  \vv{v} \\ &=3\| \vv{u}\|^{2}+7  \vv{u} \cdot  \vv{v}+2\| \vv{v}\|^{2} \\ &=3 \times 6^{2}+7 \times 2+2 \times 5^{2}=108+14+50=172 \end{aligned}$
\end{enumMet}

\subsection*{Ex 52 p 246}

\begin{enumMet}
    \item On a $ \vv{F}=\vv{F_{1}}+\vv{F_{2}}$.
    \item On a:\\
$\begin{aligned}\|\vv{F}\|^{2}=&\left\|\vv{F}_{1}+\vv{F}_{2}\right\|^{2}=\left\|\vv{F}_{1}\right\|^{2}+2 \vv{F}_{1} \cdot \vv{F}_{2}+\left\|\vv{F}_{2}\right\|^{2} \\=& 15^{2}+2 \times F_{1} \times F_{2} \times \cos \left(\vv{F}_{1}, \vv{F}_{2}\right)+13^{2} \\=& 225+2 \times 15 \times 13 \times \cos (40)+169 \\=& 394+390 \times \cos (40) \approx 692,8 \\ \text { donc }\| \vv{F}\| & \approx 26,3 \end{aligned}$
\end{enumMet}

\section*{Produit scalaire par projection orthogonale}

\subsection*{Ex 29 p 243}

\begin{enumMet}
    \item $\vv{A B} \cdot \vv{A C}=\vv{A B} \cdot \vv{A B}=c^{2}$ par projection orthogonale du point $C$ sur le segment $[A B]$.
    \item $\vv{A B} \cdot \vv{A O}=\vv{A B} \cdot \dfrac{1}{2} \vv{A B}=\dfrac{c^{2}}{2}$ par projection orthogonale du point $O$ sur le segment $[A B]$
    \item $\vv{A B} \cdot \vv{C D}=-\vv{A B} \cdot \vv{A B}=-c^{2}$
    \item $\vv{B C} \cdot \vv{A D}=\vv{B C} \cdot \vv{B C}=c^{2}$
    \item $\vv{D B} \cdot \vv{A B}=\vv{A B} \cdot \vv{A B}=c^{2}$ par projection orthogonale du point $D$ sur le segment $[A B]$
    \item $\vv{O A} \cdot \vv{A C}=-\dfrac{1}{2} \vv{A C} \cdot \vv{A C}=-\dfrac{1}{2}(\sqrt{2} c)^{2}=-c^{2}$
    \item $\vv{DB} \cdot \vv{OC}=0$ car les 2 vecteurs sont orthogonaux.
\end{enumMet}

\subsection*{Ex 39 p 244}

$\vv{A B} \cdot \vv{A C}$ avec $\vv{A B} \cdot \vv{A B}$ par projection orthogonale du point $C$ sur le segment $[A B]$\\
$\vv{A G} \cdot \vv{A F}$ avec $\vv{A G} \cdot \vv{A D}$ par projection orthogonale du point $D$ sur la droite $(A G)$\\
$\vv{A F} \cdot \vv{A B}$ avec $\vv{A H} \cdot \vv{A B}$ par projection orthogonale du point $F$ sur le segment $[A B]$\\
$\vv{A D} \cdot \vv{A F}$ avec $\vv{A D} \cdot \vv{A D}$ par projection orthogonale du point $F$ sur le segment $[A D]$

\subsection*{Ex 40 p 244}
\begin{enumMet}
    \item On projette le vecteur $\vv{A C}$ sur le vecteur $\vv{A B}$ :\\
    $\vv{A B} \cdot \vv{A C}=\vv{A B} \cdot \vv{A B}=3 \times 3=9$
    \item On projette le vecteur $\vv{B F}$ sur le vecteur $\vv{B E}$ :\\
    $\vv{B A} \cdot \vv{B F}=\vv{B A} \cdot \vv{B E}=-3 \times 2=-6$
    \item On projette le vecteur $\vv{E I}$ sur le vecteur $\vv{E B}$ :\\
    $\vv{E I} \cdot \vv{A G}=\vv{E B} \cdot \vv{A G}=-2 \times 7=-14$
    \item On projette le vecteur $\vv{G D}$ sur le vecteur $\vv{H D}$ :\\
    $\vv{C F} \cdot \vv{G D}=\vv{C F} \cdot \vv{H D}=-2 \times 7=-14$
    \item Les droites $(B C)$ et $(G H)$ sont parallèles donc on projette le vecteur $\vv{I C}$ sur le vecteur $\vv{H G}$.\\
    Pour calculer la longueur $I C$, on utilise le théorème de Thalès dans le triangle $G H D$\\
    $\dfrac{I C}{G H}=\dfrac{D C}{D H}$, donc $ I C=\dfrac{D C \times G H}{D H}=\dfrac{3 \times 3}{7}=\dfrac{9}{7}$\\
    Alors $\vv{I C} \cdot \vv{H G}=-\dfrac{9}{7} \times 3=-\dfrac{27}{7}$
    \item Pour calculer la longueur $E J$, on utilise le théorème de Thalès dans le triangle $A G D$\\
    $\dfrac{E J}{A D}=\dfrac{G E}{G A}$, donc $ E J=\dfrac{G E \times A D}{G A}=\dfrac{2 \times 3}{7}=\dfrac{6}{7}$\\
    On projette le vecteur $\vv{F A}$ sur le vecteur $\vv{F E}$ :\\
    $\vv{E J} \cdot \vv{F A}=\vv{E J} \cdot \vv{F E}=-\dfrac{6}{7} \times 3=-\dfrac{18}{7}$
\end{enumMet}

\subsection*{Ex 44 p 244}
\begin{enumMet}
    \item $A B C$ est un triangle équilatéral, donc $(\vv{B C}, \vv{B A})=\dfrac{\pi}{3}$.\\
    Alors $\vv{B C} \cdot \vv{B A}=\|\vv{B C}\| \times\|\vv{B A}\| \times \cos \left(\dfrac{\pi}{3}\right)=2 \times 2 \times \dfrac{1}{2}=2$
    \item I étant le pied de la hauteur issue de $A$ dans le triangle équilatéral $A B C$, il est également le milieu du segment $[B C],$ donc $B I=1 \mathrm{~cm}$.\\
    De plus, la symétrie du produit scalaire implique que $\vv{B A} \cdot \vv{B I}=\vv{B I} \cdot \vv{B A}$ et comme $(\vv{B I}, \vv{B A})=\dfrac{\pi}{3}$, on obtient:\\
    $\vv{B A} \cdot \vv{B I}=\vv{B I} \cdot \vv{B A}=\|\vv{B I}\| \times\|\vv{B A}\| \times \cos \left(\dfrac{\pi}{3}\right)=1 \times 2 \times \dfrac{1}{2}=1$\\
    (Ou bien $\vv{B A} \cdot \vv{B I}=\vv{B I} \cdot \vv{B I}=1^{2}=1$ en utilisant le projeté orthogonal).
    \item La droite $(A I)$ est également la bissectrice du sommet $A$ du triangle $A B C$, donc $(\vv{A I}, \vv{A C})=\dfrac{\pi}{6}$.\\
    Pour trouver la longueur du segment $[A I]$, on utilise le théorème de Pythagore dans le triangle rectangle $A I C$ :\\
    $A C^{2}=A I^{2}+C I^{2},$ donc $A I^{2}=A C^{2}-C I^{2}=2^{2}-1^{2}=4-1=3$ et enfin
    $A I=\sqrt{3} \mathrm{~cm} .$ On obtient alors:\\
    $\vv{A I} \cdot \vv{A C}=\|\vv{A I}\| \times\|\vv{A C}\| \times \cos \left(\dfrac{\pi}{6}\right)=\sqrt{3} \times 2 \times \dfrac{\sqrt{3}}{2}=3$
\end{enumMet}

%\subsection*{Ex 53 p 246}
%\begin{enumMet}
%    \item $\|\vv{A C}\|=\|\vv{B D}\|=\sqrt{5^{2}+3^{2}}=\sqrt{25+9}=\sqrt{34}$
%    \item $\vv{A C}+\vv{B D}=2 \vv{A O}+2 \vv{O D}=2 \vv{A D}$ donc $\|\vv{A C}+\vv{B D}\|=2\|\vv{A D}\|=6$
%    \item On a alors $\vv{A C} \cdot \vv{B D}=\dfrac{1}{2}\left(\|\vv{A C}+\vv{B D}\|^{2}-\|\vv{A C}\|^{2}-\|\vv{B D}\|^{2}\right)=-16$
%\end{enumMet}

\section*{Produit scalaire dans un repère orthonormée}

\subsection*{Ex 26 p 243}
\begin{enumMet}
    \item $\vv{A B} \cdot \vv{A C}=0 \times 5+(-2) \times(-1)=2$
    \item $\vv{A B} \cdot \vv{A C}=3 \times 4+1 \times(-3)=12-3=9$
    \item $\vv{A B} \cdot \vv{A C}=7 \times 5+(-4) \times 6=35-24=11$
\end{enumMet}

%\newpage
\subsection*{Ex 37 p 244}
\begin{enumMet}
    \item $ \vv{u} \cdot  \vv{v}=1 \times(-2)+0 \times 3=-2$
    \item $ \vv{u} \cdot  \vv{v}=3 \times \dfrac{5}{3}+\frac{-4}{5} \times 2=5-\dfrac{8}{5}=\frac{25-8}{5}=\dfrac{17}{5}$
    \item $ \vv{u} \cdot  \vv{v}=\sqrt{3} \times \sqrt{12}+5 \times 2=\sqrt{36}+10=6+10=16$
    \item $ \vv{u} \cdot  \vv{v}=2 x \times(-x+3)+y \times(y-1)=-2 x^{2}+6 x+y^{2}-y$
\end{enumMet}

\section*{Calcul d'angles à l'aide du produit scalaire}

\subsection*{Ex 41 p 245}
$\| \vv{u}\|=\sqrt{2^{2}+(-2)^{2}}=\sqrt{4+4}=\sqrt{8}$\\
$\| \vv{v}\|=\sqrt{3^{2}+1^{2}}=\sqrt{9+1}=\sqrt{10}$\\
$ \vv{u} \cdot  \vv{v}=2 \times 3+(-2) \times 1=6-2=4$\\
$\vv{u} \cdot \vv{v}=\|\vv{u}\| \times\|\vv{v}\| \times \cos \left(\vv{u};\vv{v}\right)=\sqrt{8}\times\sqrt{10}\times\cos \left(\vv{u};\vv{v}\right)$\\
$\cos \left(\vv{u};\vv{v}\right)=\dfrac{4}{\sqrt{8}\times\sqrt{10}}$\\
$\left(\vv{u};\vv{v}\right)\approx63,4\degree$

\subsection*{Ex 26 p 243}
Calculs d'angle
\begin{enumMet}
	\item $\vv{A B} \cdot \vv{A C}=0 \times 5+(-2) \times(-1)=2$\\
	$AB=\sqrt{0^2+(-2)^2}=2$\quad  et\quad $AC=\sqrt{5^2+(-1)^2}=\sqrt{26}$\\
	$\vv{AB} \cdot \vv{AC}=AB \times AC \times \cos \left(\widehat{BAC}\right)=2\times\sqrt{26}\times\cos \left(\widehat{BAC}\right)$\\
	$\cos \left(\widehat{BAC}\right)=\dfrac{2}{2\times\sqrt{26}}$\\
	$\widehat{BAC}\approx 78,7\degree$
	\item $\vv{A B} \cdot \vv{A C}=3 \times 4+1 \times(-3)=12-3=9$\\
	$AB=\sqrt{3^2+1^2}=\sqrt{10}$\quad  et\quad $AC=\sqrt{4^2+(-3)^2}=5$\\
	$\vv{AB} \cdot \vv{AC}=AB \times AC \times \cos \left(\widehat{BAC}\right)=\sqrt{10}\times5\times\cos \left(\widehat{BAC}\right)$\\
	$\cos \left(\widehat{BAC}\right)=\dfrac{9}{\sqrt{10}\times5}$\\
	$\widehat{BAC}\approx 55,3\degree$
	\item $\vv{A B} \cdot \vv{A C}=7 \times 5+(-4) \times 6=35-24=11$\\
	$AB=\sqrt{7^2+(-4)^2}=\sqrt{65}$\quad  et\quad $AC=\sqrt{5^2+6^2}=\sqrt{61}$\\
	$\vv{AB} \cdot \vv{AC}=AB \times AC \times \cos \left(\widehat{BAC}\right)=\sqrt{65}\times\sqrt{61}\times\cos \left(\widehat{BAC}\right)$\\
	$\cos \left(\widehat{BAC}\right)=\dfrac{11}{\sqrt{65}\times\sqrt{61}}$\\
	$\widehat{BAC}\approx 79,9 \degree$
\end{enumMet}

\section*{Exercices de synthèse}


\subsection*{Ex 43 p 245}
\begin{enumMet}
    \item $ \vv{u} \cdot  \vv{v}=(1-\sqrt{2})(1+\sqrt{2})+1 \times(-4)=1-2-4=-5$


    \item $(3  \vv{u}) \cdot  \vv{v}=3( \vv{u} \cdot  \vv{v})=3 \times(-5)=-15$
    \item $(4  \vv{u}) \cdot(-2  \vv{v})=4 \times(-2) \times( \vv{u} \cdot  \vv{v})=(-8) \times(-5)=40$
    \item \ \\
    $\begin{aligned}  \vv{u} \cdot(\sqrt{2}  \vv{u}- \vv{v}) &= \vv{u} \cdot \sqrt{2}  \vv{u}- \vv{u} \cdot  \vv{v} \\ &=\sqrt{2}( \vv{u} \cdot  \vv{u})-(-5) \\ &=\sqrt{2}\left((1-\sqrt{2})^{2}+1^{2}\right)+5 \\ &=\sqrt{2}(1-2 \sqrt{2}+2+1)+5 \\ &=\sqrt{2}(4-2 \sqrt{2})+5=4 \sqrt{2}-4+5=4 \sqrt{2}+1 \end{aligned}$
\end{enumMet}

\subsection*{Ex 44 p 245}
\begin{enumMet}
    \item 1. $A B C$ est un triangle équilatéral, donc $(\vv{B C}, \vv{B A})=\dfrac{\pi}{3}$.\\
    Alors $\vv{B C} \cdot \vv{B A}=\|\vv{B C}\| \times\|\vv{B A}\| \times \cos \left(\dfrac{\pi}{3}\right)=2 \times 2 \times \dfrac{1}{2}=2$\\
    \item I étant le pied de la hauteur issue de $A$ dans le triangle équilatéral $A B C$, il est également le milieu du segment $[B C]$, donc $B I=1 \mathrm{~cm}$.\\
De plus, la symétrie du produit scalaire implique que $\vv{B A} \cdot \vv{BI}=\vv{BI} \cdot \vv{B A}$ et comme $(\vv{B I}, \vv{B A})=\dfrac{\pi}{3}$, on obtient:\\
$\vv{B A} \cdot \vv{B I}=\vv{B I} \cdot \vv{B A}=\|\vv{B I}\| \times\|\vv{B A}\| \times \cos \left(\dfrac{\pi}{3}\right)=1 \times 2 \times \frac{1}{2}=1$.
(Ou bien $\vv{B A} \cdot \vv{B I}=\vv{B I} \cdot \vv{B I}=1^{2}=1$ en utilisant le projeté orthogonal).
\item La droite $(A I)$ est également la bissectrice du sommet $A$ du triangle $A B C$, donc $(\vv{A I}, \vv{A C})=\dfrac{\pi}{6}$. Pour trouver la longueur du segment $[A I]$, on utilise le théorème de Pythagore dans le triangle rectangle $A I C$ :
$A C^{2}=A I^{2}+C I^{2}$, donc $A I^{2}=A C^{2}-C I^{2}=2^{2}-1^{2}=4-1=3$ et enfin $A I=\sqrt{3} \mathrm{~cm}$. On obtient alors :\\
$\vv{A I} \cdot \vv{A C}=\|\vv{A I}\| \times\|\vv{A C}\| \times \cos \left(\dfrac{\pi}{6}\right)=\sqrt{3} \times 2 \times \dfrac{\sqrt{3}}{2}=3$
\end{enumMet}

\subsection*{Ex 45 p 245}
\begin{enumMet}
    \item $\vv{D A} \cdot \vv{D B}=\vv{D A} \cdot \vv{D A}=36$ par projection orthogonale du point $B$ sur le segment $[D A]$.
    \item  $\vv{D C} \cdot \vv{D F}=\vv{D C} \cdot \vv{D C}=100$ par projection orthogonale du point $F$ sur le segment $[D C]$.
    \item $\vv{D A} \cdot \vv{D C}=0$ car $A B C D$ est un rectangle.
    \item Pour déterminer le produit scalaire $\vv{D E} \cdot \vv{D F}$, on peut par exemple se placer dans le repère orthonormé $\left(A ; \dfrac{1}{10} \vv{A B}, \dfrac{1}{6} \vv{A D}\right)$.\\
    Alors, les coordonnées du point $D$ sont $(0 ; 6)$, celles de $E$ sont $(5 ; 0)$ et celles de $F$ sont $(10 ; 3)$. On calcule alors les coordonnées des vecteurs $\vv{D E}$ et $\vv{D F}$ afin d'utiliser la formule du produit scalaire dans un repère : $\vv{D E}\left(\begin{array}{c}5 \\ -6\end{array}\right)$ et $\vv{D F}\left(\begin{array}{c}10 \\ -3\end{array}\right)$.
Enfin, $\vv{D E} \cdot \vv{D F}=5 \times 10+(-6) \times(-3)=50+18=68$.
\item \textbf{Méthode 1:} Pareillement, les coordonnées du point $B$ sont $(10 ; 0)$, donc $\vv{D B}\left(\begin{array}{c}10 \\ -6\end{array}\right)$ et $\vv{D F} \cdot \vv{D B}=10 \times 10+(-3) \times(-6)=100+18=118$.\\
\textbf{Méthode 2 :} On décompose les vecteurs avec la relation de Chasles :
$$
\begin{aligned}
\vv{D F} \cdot \vv{D B} &=(\vv{D C}+\vv{C F}) \cdot(\vv{D C}+\vv{C B}) \\
&=\vv{D C} \cdot \vv{D C}+\underbrace{\vv{D C} \cdot \vv{C B}}_{0}+\underbrace{\vv{C F} \cdot \vv{D C}}_{0}+\vv{C F} \cdot \vv{C B} \\
&=10 \times 10+3 \times 6=100+18=118
\end{aligned}
$$
\end{enumMet}

\subsection*{Ex 46 p 245}

    On calcule les coordonnées des vecteurs $\vv{A B}$ et $\vv{C D}$ :\\
    $\Coord{AB}{-2-5}{7-(-3)}$, donc $\Coord{AB}{-7}{10}$ et $\Coord{CD}{-5-(-\dfrac{1}{2})}{\dfrac{3}{4}-0}$, $\Coord{CD}{-\dfrac{9}{2}}{\dfrac{3}{4}}$
Puis on calcule le produit scalaire :
$$
\vv{A B} \cdot \vv{C D}=(-7) \times\left(-\frac{9}{2}\right)+10 \times \frac{3}{4}=\frac{63}{2}+\frac{30}{4}=\frac{126}{4}+\frac{30}{4}=\frac{156}{4}=39
$$
On calcule les coordonnées des vecteurs $\vv{A C}$ et $\vv{B D}$ :\\
$\Coord{AC}{-\dfrac{1}{2}-5}{0-(-3)}$, donc $\Coord{AC}{-\dfrac{11}{2}}{3}$ et $\Coord{BD}{-5-(-2)}{\dfrac{3}{4}-7}$, $\Coord{BD}{-3}{-\dfrac{25}{4}}$.\\
Puis on calcule le produit scalaire :
$$
\vv{A C} \cdot \vv{B D}=\left(-\frac{11}{2}\right) \times(-3)+3 \times\left(-\frac{25}{4}\right)=\frac{33}{2}-\frac{75}{4}=\frac{66}{4}-\frac{75}{4}=-\frac{9}{4} \text {. }
$$
On calcule les coordonnées des vecteurs $\vv{A D}$ et $\vv{B C}$ :\\
$\Coord{AD}{-5-5}{\dfrac{3}{4}-(-3)}$, donc $\Coord{AD}{-10}{\dfrac{15}{4}}$ et $\Coord{BC}{-\dfrac{1}{2}-(-2)}{0-7}$, $\Coord{BC}{\dfrac{3}{2}}{-7}$.\\
Puis on calcule le produit scalaire :
$$\vv{A D} \cdot \vv{B C}=(-10) \times \frac{3}{2}+\frac{15}{4}\times(-7)=-\frac{30}{2}-\frac{105}{4}=-\frac{60}{4}-\frac{105}{4}=-\frac{165}{4} .$$

\subsection*{Ex 49 p 246}
On a $ \vv{u} \cdot  \vv{v}=2 \times(-1)+x \times 4=4 x-2$, donc:
\begin{enumMet}
    \item $4 x-2=2 \Longleftrightarrow x=1$
    \item $4 x-2=-5 \Longleftrightarrow x=-\dfrac{3}{4}$
    \item $4 x-2=\dfrac{7}{3} \Longleftrightarrow x=\dfrac{13}{12}$
    \item $4 x-2=\sqrt{8} \Longleftrightarrow x=\dfrac{2+2 \sqrt{2}}{4}=\dfrac{1+\sqrt{2}}{2}$
\end{enumMet}


\subsection*{Ex 54 p 246}
\begin{enumMet}
    \item On calcule le produit scalaire de ces deux vecteurs : $ \vv{u} \cdot  \vv{v}=-18+2 x$. Les vecteurs sont orthogonaux si et seulement si leur produit scalaire est nul, ce qui implique que $-18+2 x=0 .$ On résout cette équation $: x=9$.\\
     Les vecteurs sont orthogonaux pour $x=9$
     \item $ \vv{u} \cdot  \vv{v}=-3(x-1)+4 x=-3 x+3+4 x=x+3$. Les vecteurs sont orthogonaux si et seulement si leur produit scalaire est nul, ce qui implique que $x=-3$.\\
     Les vecteurs sont orthogonaux pour $x=-3$.
     \item On calcule le produit scalaire de ces deux vecteurs : $ \vv{u} \cdot  \vv{v}=3 x-16$. Les vecteurs sont orthogonaux si et seulement si leur produit scalaire est nul, ce qui implique que $3 x-16=0$.\\
     Les vecteurs sont orthogonaux pour $x=\dfrac{16}{3}$
     \item $ \vv{u} \cdot  \vv{v}=x^{2}-16$. Les vecteurs sont orthogonaux si et seulement si leur produit scalaire est nul, ce qui implique que $x^{2}=16$.\\
     Les vecteurs sont orthogonaux pour $x=-4$ ou $x=4$
 \end{enumMet}
 \subsection*{Ex 55 p 246}
 \begin{enumMet}
     \item $ \vv{u} \cdot  \vv{v}=-8 t+2 t=-6 t$.\\
     Les vecteurs sont orthogonaux si et seulement si leur produit scalaire est nul, ce qui implique que $-6 t=0$.\\
     Les vecteurs sont orthogonaux pour $t=0$. 
     \item $ \vv{u} \cdot  \vv{v}=t(t+6)-2\left(2 t-\frac{1}{2}\right)=t^{2}+6 t-4 t+1=t^{2}+2 t+1=(t+1)^{2}$\\
     Les vecteurs sont orthogonaux si et seulement si leur produit scalaire est nul, ce qui implique que $(t+1)^{2}=0,$ donc $t+1=0,$ donc $t=-1$\\
     Les vecteurs sont orthogonaux pour $t=-1$
     \item \ \\
     $$
 \begin{aligned}
  \vv{u} \cdot  \vv{v} &=\sqrt{3} t \times(\sqrt{5} t+1)+2 t \times \sqrt{12} \\
 &=\sqrt{15} t^{2}+\sqrt{3} t+2 \sqrt{12} t=\sqrt{15} t^{2}+(\sqrt{3}+4 \sqrt{3}) t \\
 &=\sqrt{15} t^{2}+5 \sqrt{3} t=\sqrt{15}\left(t^{2}+\sqrt{5} t\right)=\sqrt{15} t(t+\sqrt{5})
 \end{aligned}
 $$
 Les vecteurs sont orthogonaux si et seulement si leur produit scalaire est nul, ce qui implique que $t(t+\sqrt{5})=0,$ alors les solutions de l'équation produit nul sont 0 et $-\sqrt{5}$.\\
 Les vecteurs sont orthogonaux pour deux valeurs de $t: 0$ et $-\sqrt{5}$.
 \end{enumMet}

\subsection*{Ex 56 p 247}
\begin{enumMet}
    \item On calcule les coordonnées des vecteurs $\vv{A B}, \vv{B C}$ et $\vv{C A}:$\\
    $\Coord{AB}{-4-(-10)}{1-4}$, donc $\Coord{AB}{6}{-3}$;\\
    $\Coord{BC}{-1-(-4)}{7-1}$, donc $\Coord{BC}{3}{6}$;\\
    $\Coord{CA}{-10-(-1)}{4-7}$, donc $\Coord{CA}{-9}{-3}$;\\
    On calcule les trois produits scalaires des côtés du triangle :\\
    $\vv{A B} \cdot \vv{B C}=6 \times 3+(-3) \times 6=0$\\
    $\vv{B C} \cdot \vv{C A}=3 \times(-9)+6 \times(-3)=-27-18=-45$\\
    $\vv{C A} \cdot \vv{A B}=(-9) \times 6+(-3) \times(-3)=-54+9=-45$\\
    On en déduit que le triangle $A B C$ est rectangle en $B$.
    \item \textbf{Méthode 1} : Les diagonales du rectangle $A B C D$ sont $[A C]$ et $[B D]$.\\
    Comme un rectangle est un parallélogramme, les diagonales ont le même milieu $O$ de coordonnées:\\
    $\left(\dfrac{x_{A}+x_{C}}{2} ; \dfrac{y_{A}+y_{C}}{2}\right)$ donc $ O\left(\dfrac{-10+(-1)}{2} ; \dfrac{4+7}{2}\right),$ donc $O\left(-\dfrac{11}{2} ; \dfrac{11}{2}\right)$.\\ Or on sait que\\
    $O\left(\dfrac{x_{B}+x_{D}}{2} ; \dfrac{y_{B}+y_{D}}{2}\right)$, donc $ O\left(\dfrac{-4+x_{D}}{2} ; \dfrac{1+y_{D}}{2}\right)$ Alors $ \dfrac{-4+x_{D}}{2}=-\dfrac{11}{2}$ et
    $\dfrac{1+y_{D}}{2}=\dfrac{11}{2},$ donc $x_{D}=-7$ et $y_{D}=10$\\
    Les coordonnées du point $D$ tel que le quadrilatère $A B C D$ soit un rectangle sont donc $D(-7 ; 10)$\\
    \ \\
    \textbf{Méthode 2} : Comme $A B C D$ est un parallélogramme, $\vv{A B}=\vv{D C}$, or ${A B}\left(\begin{array}{c}6 \\ -3\end{array}\right)$ donc $\vv{D C}\left(\begin{array}{c}6 \\ -3\end{array}\right).$\\
    Comme $x_{\vv{D C}}=x_{C}-x_{D}=-1-x_{D \text { et }} y_{\vv{D C}}=y_{C}-y_{D}=7-y_{D}$ on en déduit que\\
    $x_{D}=-1-6=-7$ et $y_{D}=7-(-3)=10$.\\
    Les coordonnées du point $D$ tel que le quadrilatère $A B C D$ soit un rectangle sont donc $D(-7 ; 10)$
\end{enumMet}

\subsection*{Ex 57 p 247}
On va calculer les produits scalaires $\vv{A B} \cdot \vv{C D}, \vv{A B} \cdot \vv{C E}$ et $\vv{A B} \cdot \vv{C F}$\\
On a $A B\left(\begin{array}{c}2-(-4) \\ -3-(-2)\end{array}\right)$ donc $ \vv{A B}\left(\begin{array}{c}6 \\ -1\end{array}\right),$\\
$\vv{C D}\left(\begin{array}{c}0-1 \\ -2-4\end{array}\right)$ donc $ \vv{C D}\left(\begin{array}{c}-1 \\ -6\end{array}\right)$\\
$\vv{C E}\left(\begin{array}{c}\dfrac{1}{2}-1 \\ 1-4\end{array}\right)$, donc $ \vv{C E}\left(\begin{array}{r}-\dfrac{1}{2} \\ -3\end{array}\right)$ et \\
$\vv{C F}\left(\begin{array}{c}2-1 \\ 7-4\end{array}\right)$ donc $ \vv{C F}\left(\begin{array}{l}1 \\ 3\end{array}\right)$\\

$\vv{A B} \cdot \vv{C D}=6 \times(-1)+(-1) \times(-6)=-6+6=0,$ donc les vecteurs $\vv{A B}$ et $\vv{C D} \operatorname{sont}$
orthogonaux, donc $D$ est sur la hauteur issue de $C$ dans le triangle $A B C$.\\ $\vv{A B} \cdot \vv{C E}=6 \times\left(-\frac{1}{2}\right)+(-1) \times(-3)=-3+3=0,$ donc les vecteurs $\vv{A B}$ et $\vv{C E}$ sont
orthogonaux, donc $E$ est sur la hauteur issue de $C$ dans le triangle $A B C$.\\ $\vv{A B} \cdot \vv{C F}=6 \times 1+(-1) \times 3=6-3=3 \neq 0,$ donc les vecteurs $\vv{A B}$ et $\vv{C F}$ ne sont pas orthogonaux, donc $F$ n'est pas sur la hauteur issue de $C$ dans le triangle $A B C$.




\end{document}

